\section{A finite recurrence for the algebraic calculations}
\label{app:recurrence}

This appendix specifies the finite arithmetic behind the tables in
Section~\ref{sec:finite}.  Its input is the integral Brown--Peterson Hopf
algebroid, and its operations are polynomial multiplication, coefficient
extraction, and elimination over $\F_3$.  In particular, neither an Ext
chart nor a choice of unlisted Adams--Novikov differentials is an input.
The distinction between a coboundary in the algebraic complex and a
topological spectral-sequence differential will be maintained throughout:
we write $\partial$ for a resolution differential, $d_C$ for its cochain
differential, and $L$ for the normalized possible $d_5$ of
Section~\ref{sec:finite}.

\subsection{Integral structure constants and the algebra of operators}

Divide internal degrees by four, and put
\[
 w_i=\frac{3^i-1}{2},\qquad
 |Q|=\sum_{i\geq1}Q_iw_i,\qquad
 B=\F_3[v_1,v_2,\ldots],\qquad |v_i|=|t_i|=w_i.
\]
All multi-indices have nonnegative integer entries and finite support.
The notation $[t^T]$ means coefficient extraction, leaving the coefficient
polynomial in the $v_i$.  The same convention applies to several successive
sets of $t$ variables.

Here is a recursive specification of all structure constants used below.
Work first in $\Q[v_1,v_2,\ldots,t_1,t_2,\ldots]$, put $l_0=t_0=1$, and
successively calculate
\begin{align}
 3l_n&=\sum_{j=0}^{n-1}l_jv_{n-j}^{3^j},
 &\rho_n&=\sum_{j=0}^{n}l_jt_{n-j}^{3^j},\label{alg:logarithms}\\
 R_0&=1,
 &R_n&=3\rho_n-\sum_{j=1}^{n-1}\rho_jR_{n-j}^{3^j}.
 \label{alg:right-unit}
\end{align}
Then $R_n=\eta_R(v_n)$ and $\rho_n=\eta_R(l_n)$.  Introduce a second set
$s_i$ of strict-isomorphism coordinates, with $s_0=1$.  With $D_0=1$, define
\begin{equation}\label{alg:coproduct}
 D_n=\sum_{j+k+h=n}l_jt_k^{3^j}s_h^{3^{j+k}}
          -\sum_{j=1}^{n}l_jD_{n-j}^{3^j}.
\end{equation}
Thus $D_n=\Delta(t_n)$, with $t_i=t_i\otimes1$ and
$s_i=1\otimes t_i$.  These are the standard Hazewinkel coordinates for
the integral $\BP$ Hopf algebroid; in particular the rational expressions
for $R_n,D_n$ are integral at three
\cite[Theorem~A2.1.27, Equation~(A2.2.1), and Theorem~A2.2.3]{Ravenel}.  This use of the
standard presentation is the only structure theorem needed for the
arithmetic.  Formula~\eqref{alg:coproduct} is equivalently the identity
obtained by composing two strict isomorphisms and comparing logarithms.
Coefficients in a tensor product are always written in the leftmost
factor.  Moving a coefficient from a later factor to this form uses
$\eta_R$, not the left unit.

For ordinary calculations reduce $R_n,D_n$ modulo three after performing
the integral divisions.  For a Bockstein, retain their values modulo nine
until after division by three.  Only $1\leq i\leq4$ can occur at any of
our cutoffs, because $w_5=121>93$.  For example,
\begin{align*}
 \eta_R(v_1)&=v_1,&
 \eta_R(v_2)&=v_2+v_1t_1^3-v_1^3t_1\quad\pmod3,\\
 \eta_R(v_1)&=v_1+3t_1,&
 \eta_R(v_2)&=v_2+3t_2+5v_1^3t_1+v_1t_1^3\quad\pmod9.
\end{align*}
Exact rational polynomial arithmetic followed by integral reduction is
one way to implement these instructions; no division in $\F_3$ is intended.

Let $P_T$ be the operation dual to $t^T$, and let $P_0=1$.  The left
$B$-module of operations has basis $P_T$, with multiplication and action
on $B$ prescribed by
\begin{align}
 P_Sv^QP_T
 &=\sum_{I\leq S}\ \sum_U
       [t^I]\eta_R(v^Q)\,
       [t^{S-I}s^T]\Delta(t^U)\,P_U,
       \label{alg:operator-product}\\
 P_S(v^Q)&=[t^S]\eta_R(v^Q).
       \label{alg:operator-action}
\end{align}
Denote this algebra by $\mathcal U$.  Its products are understood in the
degreewise completion specified below.  Associativity and the module
identity for~\eqref{alg:operator-action} follow respectively by extracting
coefficients from coassociativity and from the compatibility of
$\Delta$ with $\eta_R$.  Alternatively, both identities follow directly
by evaluating a composite of three strict isomorphisms in either order.
The coefficient subalgebra $B\subset\mathcal U$ need not be central.
The augmentation of its left module resolution is the $\mathcal U$-linear
map
\begin{equation}\label{alg:augmentation}
 \varepsilon:\mathcal U\longrightarrow B,
 \qquad \varepsilon(u)=u(1),\qquad
 \varepsilon(v^QP_T)=
 \begin{cases}v^Q&T=0,\\0&T\ne0.\end{cases}
\end{equation}
It is not necessary, or correct in general, to treat this augmentation
as an algebra homomorphism to $B$.

\subsection{The additive complex and its contraction}

The use of an additive associated graded and a smaller resolution belongs
to the filtered bar/cobar method developed by May
\cite[Sections~4, 6--7]{MayHopfCohomology}; the contraction, coefficient
corrections, and cutoff comparison needed here are given below.

A \emph{digit} is a pair $\nu=(i,j)$ with $i\geq1$, $j\geq0$.
Put $m_\nu=3^j\mathbf e_i$, where $\mathbf e_i$ is the $i$th unit
multi-index.  Order digits by decreasing $i$, then decreasing $j$.
At cutoff $N$ retain only digits with $|m_\nu|\leq N$.
A generator $g=(e_\nu)$ is a tuple of nonnegative integers; set
\begin{equation}\label{alg:generators}
 s(g)=\sum_\nu e_\nu,\qquad
 C(g)=\sum_\nu\left(3\left\lfloor\frac{e_\nu}{2}\right\rfloor
                         +(e_\nu\bmod2)\right)m_\nu,
 \qquad w(g)=|C(g)|.
\end{equation}
Write $g_{-\nu}$ and $g_{+\nu}$ for decrementing or incrementing
$e_\nu$.  The free module $G_s$ has these generators with $s(g)=s$.
The generator with all entries zero is denoted $g_0$.

Temporarily replace the Hopf algebroid by the additive one, in which
the $t_i$ are primitive and all coefficients are invariant.  Its dual
multiplication is
\[
 P_SP_T=\prod_i\binom{S_i+T_i}{S_i}P_{S+T}.
\]
Lucas's formula, or the coefficient identity
$(1+z)^{3^j}=1+z^{3^j}$ in characteristic three, splits this algebra into
the tensor product of the algebras $\F_3[x_\nu]/(x_\nu^3)$, where
$x_\nu=P_{m_\nu}$ and $x_\nu^2=-P_{2m_\nu}$.  The periodic resolution
for one such factor alternates multiplication by $x_\nu$ and
$x_\nu^2$.  Its tensor product has differential
\begin{equation}\label{alg:additive-differential}
 \partial^0g=\sum_{\nu:e_\nu>0}
 (-1)^{\sum_{\mu<\nu}e_\mu}
 \begin{cases}
 P_{m_\nu}g_{-\nu},& e_\nu\text{ odd},\\
 -P_{2m_\nu}g_{-\nu},& e_\nu\text{ even}.
 \end{cases}
\end{equation}

An explicit $B$-linear contraction, denoted $H$, will be useful;
this operator is distinct from the homology $H(D,d_5)$ used in
Section~\ref{sec:finite}.  In a term
$v^QP_Tg$, let $j_0\in\{0,1,2\}$ be the base-three digit of $T_i$
at the position $\nu=(i,j)$.  Select the first digit for which either
$j_0\ne0$ or $e_\nu\ne0$.  For this digit put
$h'=1+(e_\nu\bmod2)$, and define
\begin{equation}\label{alg:contraction}
 H(v^QP_Tg)=
 \begin{cases}
 j_0!\,v^QP_{T-h'm_\nu}g_{+\nu},&j_0\geq h',\\
 0,&j_0<h',
 \end{cases}
\end{equation}
with value zero if there is no active digit.

\begin{lemma}\label{alg:additive-contraction}
If $\iota:B\to G_0$ sends $b$ to $bg_0$, then
\[
 \partial^0H+H\partial^0=1-\iota\varepsilon.
\]
The map $H$ preserves the coefficient $v^Q$, the number
$|T|+w(g)$, and the tuple $T+C(g)$.
\end{lemma}

\begin{proof}
For a single digit, let $e$ be the resolution degree and let
$j_0=0,1,2$ be the operator digit.  If $e$ is even, the next boundary
multiplies by $P_m$.  On the two terms for which $H$ is nonzero its
coefficient is $j_0!j_0=1$ in $\F_3$.  If $e$ is odd, the only nonzero
case has $j_0=2$; the next boundary is $-P_{2m}$ and the combined
coefficient is $-2=1$.  On the other terms $H\partial^0$ supplies the
identity, except at $e=j_0=0$, which is the augmentation summand.

In the tensor product all factors before the first active one have
operator digit and resolution degree zero.  Hence no additional tensor
sign is required in~\eqref{alg:contraction}.  Terms differentiating a
later factor cancel between $\partial^0H$ and $H\partial^0$, because
$H$ changes the degree of the selected factor by one.  Finally, the
increase of $C(g)$ on incrementing this factor is exactly $h'm_\nu$.
This proves all the assertions.
\end{proof}

\subsection{Triangular deformation, finite cutoffs, and exactness}

We now deform $\partial^0$ to a differential for the actual operator
algebra. We order terms so that correcting the leading error leaves only
larger terms; consequently only finitely many corrections affect a fixed
weight.

The passage from the additive contraction to a deformed resolution is
a homological-perturbation construction; see \cite{Lambe1991} for this
method of constructing resolutions and
\cite[Section~4.8.3, Theorem~51]{RubioSergeraert2013} for the basic
perturbation lemma.  The completed, semilinear Brown--Peterson
implementation and its finite-cutoff arguments are supplied below,
rather than taken from those references.

Give a term $v^QP_Tg$ internal degree
\begin{equation}\label{alg:dual-degree}
 |Q|-|T|-w(g)
\end{equation}
and \emph{pure weight} $|T|+w(g)$.  Within a fixed homogeneous degree,
complete by allowing this pure weight to tend to infinity.  There are
only finitely many terms at each fixed pure weight.  Order terms first
by increasing pure weight and then by increasing tuple $T+C(g)$,
using lexicographic order in which the largest index is most significant.
Equal terms, including their coefficient monomials, are always combined
before a leading part is selected.  At pure weight at most $N$ the set
of order values is finite.

Two elementary triangularity properties of the structure constants
are needed.  Modulo three the ideal $(v_1,v_2,\ldots)$ is invariant
under $\eta_R$.  To check this without interchanging division and
reduction, let $a_i=R_i|_{v_1=v_2=\cdots=0}$ in the integral
polynomial ring.  Equations~\eqref{alg:logarithms} and
\eqref{alg:right-unit} give
\[
 a_i=3t_i-\sum_{j=1}^{i-1}t_j a_{i-j}^{3^j}.
\]
Induction gives $a_i\in3\Z_{(3)}[t_1,t_2,\ldots]$, proving the
assertion.  Multiplicativity then preserves every power of this ideal.
Also, after setting the $v_i$ to zero,
\begin{equation}\label{alg:coefficient-free-coproduct}
 \Delta(t_i)=\sum_{k=0}^{i}t_k\otimes t_{i-k}^{3^k}.
\end{equation}
Its two end terms are the primitive terms.  Every other term replaces
one occurrence of $t_i$ by variables of strictly smaller index.

It follows that a product of pure operators has the additive product
as its leading term.  Every further term with a positive-weight coefficient
has larger pure weight, by homogeneity.  Every coefficient-free further
term has larger output tuple: a nonprimitive term in the coproduct
lowers the input tuple relative to the tuple of the element whose
coproduct is being taken.  These assertions continue to hold when
correction terms are multiplied from the left.  A preexisting left
coefficient is unchanged, and a coefficient in the positive-degree
coefficient ideal remains in that ideal when moved past a left operator, by invariance
of the coefficient ideal.  For a coefficient-free correction, adding
the left operator tuple preserves a strict lexicographic inequality.
Consequently any degree-zero differential whose values on generators
have leading part~\eqref{alg:additive-differential} acts on arbitrary
module terms with leading differential $\partial^0$.

Here and below, truncation means deletion of all output terms of pure
weight greater than the stated cutoff.  Suppose $\partial_{r+1}$ is
already defined and is exact onto cycles in degree $r$.  Define
$\operatorname{Ant}_r(z;N)$ for a cycle $z\in G_r$ by the following
successive correction procedure.  At $r=0$, ``cycle'' means an element
of $\ker\varepsilon$.
\begin{equation}\label{alg:antidifferential}
\begin{array}{l}
 u\leftarrow0;\quad e\leftarrow z;\\
 \text{while }e\ne0:\\
 \quad p\leftarrow\text{the sum of terms at the least order occurring in }e;\\
 \quad z_1\leftarrow H(p);\\
 \quad u\leftarrow u+z_1;\qquad
 e\leftarrow e-\partial_{r+1}(z_1)\quad\text{(truncate at }N\text{)};\\
 \text{return }u.
\end{array}
\end{equation}
The reference to exactness in this description is justified by the
induction in the following proposition; an implementation only needs
the already constructed differential.

\begin{proposition}\label{alg:resolution-correct}
Starting with $\partial_1=\partial^0_1$ and augmentation
\eqref{alg:augmentation}, the recursion
\begin{equation}\label{alg:deformation}
 \partial_sg=p-
   \operatorname{Ant}_{s-2}(\partial_{s-1}p;N),
 \qquad p=\partial^0_sg,\qquad s\geq2,
\end{equation}
defines compatible finite truncations of a completed free resolution
of the $\mathcal U$-module $B$.  Every call to
\eqref{alg:antidifferential} terminates.  Cochains of internal weight
$d\leq N$ in this resolution calculate the ordinary normalized-cobar
cohomology $D_{s,d}$.
\end{proposition}

\begin{proof}
First consider a cycle $e$ whose leading part is $p$.  The leading part
of its boundary is $\partial^0p$, so $\partial^0p=0$.  In degree zero
there is no $P_0g_0$ term, because $\varepsilon(e)=0$ and distinct
coefficient monomials are independent.  Lemma~\ref{alg:additive-contraction}
therefore gives $\partial^0H(p)=p$ in every relevant degree.  Subtracting
$\partial H(p)$ removes precisely the leading part and introduces only
larger order values.  It preserves the cycle condition.  Since only
finitely many order values occur below $N$, the procedure terminates
and returns a preimage of $e$.

For the induction on $s$, $\varepsilon\partial_1=0$.  The preceding
argument with $r=0$ proves exactness at $G_0$.  Assume the differentials
through $\partial_{s-1}$ have been constructed, their successive
composites vanish, and cycles in $G_{s-2}$ can be lifted.  The error
$\partial_{s-1}p$ in~\eqref{alg:deformation} is such a cycle.  Its
additive leading part vanishes because $(\partial^0)^2=0$, so all
corrections have order strictly larger than that of $g$.
Equation~\eqref{alg:deformation} gives
$\partial_{s-1}\partial_sg=0$ and preserves the required leading part.
The first paragraph now lifts every cycle in $G_{s-1}$.  This completes
the induction and proves exactness at each finite cutoff.

We give the completion and comparison details to ensure that finite
truncation has not changed the invariant being calculated.  From
\eqref{alg:operator-product}, an output operator $P_U$ in
$P_Sv^QP_T$ satisfies
\begin{equation}\label{alg:right-weight-bound}
 |U|\geq |T|.
\end{equation}
To see this, take a contributing coefficient of
$\Delta(t^U)$.  If its coefficient has weight $c\geq0$, homogeneity
gives $|U|=|S-I|+|T|+c$.  Thus an output expansion on a free generator
may be truncated before acting on it from the left: discarded output
terms cannot return below the cutoff.  Input truncation is justified
separately.  Degree-zero differentials preserve a coefficient already
on the far left, and their generator values have nonnegative new
coefficient weight, so they do not decrease pure weight.  The
contraction preserves it.  The constructions at successive cutoffs
therefore agree under truncation.  Taking their compatible limits
gives exact lifting in the completed modules, rather than relying on
an exactness theorem for an infinite dual.

For comparison use the normalized relative bar resolution of $B$ over
$\mathcal U$ relative to $B$.  It is left free on the words
$[P_{T_1}|\cdots|P_{T_s}]$, $T_i\ne0$.  This description remains valid
for noncentral $B$: move coefficient factors left by the balancing
relation and~\eqref{alg:operator-product}.  The differential multiplies
successive entries with the alternating bar signs, including the first
left multiplication and the last action on $1\in B$; an entry in $B$
at a normalized position is discarded.  For example its degree-two
boundary is
\[
 P_S[P_T]-[P_SP_T]+[P_S]P_T(1),
\]
with the middle term expanded and normalized.  A $B$-linear contracting
homotopy inserts the left $P_T$ as the first word entry when $T\ne0$
and is zero when $T=0$, while keeping its coefficient on the far left.
The usual adjacent terms in the two composites cancel; the sole
uncancelled term is the original word, with the augmentation term
removed in degree zero.  This verifies the contraction directly.
It preserves pure weight, where a word generator has weight
$\sum|T_i|$.

Construct a chain map from $G$ to this bar resolution by sending $g_0$
to the empty word and, on each further generator, applying the bar
contraction to the image of its boundary.  Construct the reverse map
using~\eqref{alg:antidifferential}.  Their composites are homotopic to
the identity: after the homotopy has been defined in lower degrees,
its next required value is a cycle, and the corresponding contraction
or antidifferential supplies a preimage.  This also proves independence
of all choices of lifts.  Every map and homotopy has the cutoff
compatibility just proved.

Finally a cochain of weight $d$ takes a generator of weight $w(g)$ to
a polynomial of weight $d-w(g)$.  Evaluation of $v^QP_Tg$ uses
$P_T$ on that polynomial, which is zero if $|T|+w(g)>d$.
Only finitely many terms of any comparison or homotopy can consequently
be seen by such a cochain.  Cochains on the relative bar identify with
the ordinary cobar by coefficient extraction: multiplying adjacent
operators is dual to inserting $\Delta$, and balancing coefficients
is dual to applying $\eta_R$.  Thus the comparison is a cochain
homotopy equivalence in each weight $d\leq N$, as asserted.
\end{proof}

The same argument applies to a free-generator map of internal degree
$h$: if the final cochain weight is at most $N$, its output expansion
needs only pure weight at most $N-h$.  This is the shifted cutoff in
the product construction below.  A common resolution through weight
$93$ and filtration $33$ suffices for all the tables.  Smaller cutoffs
$60$ and $92$ suffice for their respective affine systems; the cubic
check alone needs weight $81$ and filtration $7$.

\subsection{Cochain coordinates, Bocksteins, and products}

At bidegree $(s,d)$ the cochain basis is
\begin{equation}\label{alg:cochain-basis}
 \mathcal C_{s,d}=\{(Q,g):s(g)=s,\ |Q|+w(g)=d\}.
\end{equation}
The basis vector $(Q,g)$ has value $v^Q$ on $g$ and zero on the other
generators.  If
$\partial_{s+1}g'=\sum a_{U,T,g}v^UP_Tg$, the matrix entry of $d_C$
in its column $(Q,g)$ and row $(Q',g')$ is
\begin{equation}\label{alg:cochain-matrix}
 \sum_{U,T}a_{U,T,g}
       [v^{Q'-U}t^T]\eta_R(v^Q).
\end{equation}
A coefficient with a negative exponent is zero.  This formula is also
a direct specification of every scalar in the cochain matrices.
The dimensions may be checked before any matrix calculation by
\begin{equation}\label{alg:dimension-series}
 \sum_{s,d}\#\mathcal C_{s,d}\,z^st^d
 =\prod_{i\geq1}\frac1{1-t^{w_i}}
    \prod_{\nu}\frac{1+zt^{|m_\nu|}}
                       {1-z^2t^{3|m_\nu|}}.
\end{equation}

Write $M_{s,d}$ for the matrix in~\eqref{alg:cochain-matrix}.
Choose independent incoming columns $B_{s,d}$ of $M_{s-1,d}$.
Extend them, using a kernel basis of $M_{s,d}$, to a basis
$[B_{s,d}\mid V_{s,d}]$ of $\ker M_{s,d}$.  The columns of $V_{s,d}$
are representatives of the chosen basis of $D_{s,d}$.  For a cycle
$z$, solve
\begin{equation}\label{alg:quotient-coordinates}
 z=B_{s,d}\lambda+V_{s,d}\mu;
 \qquad [z]\longleftrightarrow\mu.
\end{equation}
This avoids making names for additional Ext generators.

For completeness, the required comparison of named cobar cocycles to
these coordinates is especially short.  Denote this comparison by
$\operatorname{get}$.  In filtrations zero and one,
\[
 \operatorname{get}(q)(g_0)=q,\qquad
 \operatorname{get}(x)(g_\nu)=[t^{m_\nu}]x.
\]
Here $g_\nu$ has its single nonzero entry $e_\nu=1$.  For a normalized
two-cocycle $x$, set
\begin{equation}\label{alg:get-two}
 \operatorname{get}(x)(g)
   =\sum_{v^QP_Tg_\nu\text{ in }\partial_2g}
       a_{Q,T,\nu}\,v^Q[t^T\otimes t^{m_\nu}]x.
\end{equation}
This is the degree-two bar comparison in the proof of
Proposition~\ref{alg:resolution-correct}: its contraction inserts the
first operator entry.  The term with $T=0$ has zero coefficient in a
normalized two-cocycle.  Applying~\eqref{alg:quotient-coordinates} gives
the class of each named element.

Use the names of Section~\ref{sec:finite}; explicitly, with $v=v_1$,
\begin{align*}
 a&=[t_1],& h&=\delta_1(v_2),& r&=\delta_3(v_2^3),\\
 s'&=\delta_1(v_2^2),&
 y&=\delta_9(v_2^9-v^8v_2^7),&
 t_i'&=\delta_1(v_2^i),
\end{align*}
where $\delta_m(q)=(\eta_R(q)-q)/v^m$ is calculated after reduction
modulo three.  For these polynomials divisibility follows directly from
$\eta_R(v_2)-v_2=v(t_1^3-v^2t_1)$.  Their cocycle identities follow by
applying the cochain differential to a coefficient coboundary and
using invariance of $v$.  Notice the exact identity
\begin{equation}\label{alg:y-convention}
 \delta_9(v_2^9)=y+t_7'.
\end{equation}

The notation $B(x)$ below denotes a Bockstein cochain; the unadorned
$B$ still denotes the coefficient ring.
If $x$ is a normalized one-cocycle modulo three, lift its coefficients
integrally and set
\begin{equation}\label{alg:bockstein}
 B(x)=\frac{\Delta\widetilde x-\widetilde x\otimes1
                           -1\otimes\widetilde x}{3}\pmod3.
\end{equation}
The numerator is divisible by three because $x$ is a cocycle.  The
last tensor uses the right unit on every coefficient of $\widetilde x$.
The identity $d_C^2=0$ proves that this is a two-cocycle, and changing
the lift changes it by a coboundary.  With the convention
$d_Cx=1\otimes x-\Delta x+x\otimes1$, it is minus the mod-three
reduction of the connecting homomorphism for multiplication by three.
In particular
\begin{equation}\label{alg:b-sign-check}
 b=B(h)=t_1^2\otimes t_1+t_1\otimes t_1^2
                           -v\,t_1\otimes t_1.
\end{equation}
Set $c'=B(r)$ and $P=aB(t_4')$.  Equations
\eqref{alg:bockstein} and~\eqref{alg:get-two} also specify $B(y)$
without an integral resolution: all integral arithmetic takes place in
the ordinary cobar before the comparison.

Here is a complete product rule on the smaller resolution.  If a
cocycle $p$ has bidegree $(j,h)$, construct $\mathcal U$-linear maps
$F_s:G_{s+j}\to G_s$ of internal degree $h$ by
\begin{align}
 F_0(g)&=p(g)P_0g_0,\label{alg:lift-base}\\
 F_s(g)&=(-1)^j\operatorname{Ant}_{s-1}
       \bigl(F_{s-1}(\partial_{s+j}g);N-h\bigr),\qquad s\geq1.
       \label{alg:lift-recursion}
\end{align}
Whenever a map is extended to a term $v^QP_Tg$, use
\eqref{alg:operator-product}, multiplying the preexisting $v^Q$ on the
left last.  At $s=1$ the required augmentation condition is exactly the
cocycle condition for $p$.  In higher degrees the required cycle
condition follows inductively from
\begin{equation}\label{alg:signed-lift-identity}
 \partial F_s=(-1)^jF_{s-1}\partial.
\end{equation}
Consequently, for a cocycle $x$ of bidegree $(f,d)$, the cochain
\begin{equation}\label{alg:product}
 p\star x=xF_f
\end{equation}
is a cocycle of bidegree $(f+j,d+h)$.  Reduce it using
\eqref{alg:quotient-coordinates}.  For a slot $T=(f,d)$, write
$\operatorname{mult}(p,T)$ for the matrix whose columns are these
coordinates as $x$ ranges over the columns of $V_{f,d}$.
Iterated multiplication, in particular by $b^n$, is composition of
these matrices in the indicated successive slots.

There is a useful exact shortcut for a one-cocycle $\delta_m(q)$.
Let $F_q(g)=qg$, extended $\mathcal U$-linearly.  Since $v$ is invariant,
it is central in $\mathcal U$ modulo three.  Thus
\begin{equation}\label{alg:delta-lift}
 F_{\delta_m(q)}=\frac{F_q\partial-\partial F_q}{v^m}
\end{equation}
is well defined, and its value is obtained by replacing each term
$v^QP_Tg'$ in $\partial g$ by
\begin{equation}\label{alg:delta-lift-terms}
 \sum_{0<I\leq T}v^Q[t^I]\delta_m(q)\,P_{T-I}g'.
\end{equation}
Indeed this is precisely~\eqref{alg:operator-product} for the
commutator with $q$, with its $I=0$ term cancelled.  Its augmentation
in filtration one is the required digit coefficient of $\delta_m(q)$,
and $\partial F_{\delta_m(q)}=-F_{\delta_m(q)}\partial$ follows by
expanding the commutator and using $\partial^2=0$.  The shortcut hence
computes the same products as~\eqref{alg:lift-recursion}.

\begin{lemma}\label{alg:product-conventions}
The rule~\eqref{alg:product} calculates signed Yoneda products, with
the first factor lifted first.  It is independent on cohomology of all
choices of lifts and representatives.  It agrees with the ordinary Ext
product with the corresponding ordering convention.  If the cobar
comparison instead uses the ordinary opposite ordering of the two
factors, graded commutativity changes the product formula by
$(-1)^{jf}$ for factors of filtrations $j$ and $f$.
This alternative does not alter the use of the affine
systems or of their kernel and image bounds.
\end{lemma}

\begin{proof}
The difference of two lifts with the same augmentation is null
homotopic: in each successive degree its uncancelled difference is a
cycle, which is lifted by~\eqref{alg:antidifferential}.  The same
argument applied to a coboundary representative proves independence
of that representative.  Composition of such lifts is the defining
Yoneda composition.

One may also compare on the relative bar.  An unsigned lift applies
the cocycle to the last block of entries and retains the first block;
move the resulting coefficient left by the balancing relation and
discard scalar entries at normalized positions.  The interior terms
cancel in its chain identity, and the terms crossing the division
between blocks cancel by the cocycle equation.  Multiplying its
component in degree $s$ by $(-1)^{js}$ gives
\eqref{alg:signed-lift-identity}.  Evaluation by the other cocycle is
the bar, and hence cobar, cup product, with the stated ordering sign.

To check the final assertion explicitly, suppose an odd-filtration
differential $D$ is a left derivation for a product juxtaposed as $xy$,
and change to $x\star y=(-1)^{fg}xy$, where $x,y$ have filtrations
$f,g$.  Then
\[
 D(x\star y)=(-1)^gD(x)\star y+x\star D(y).
\]
The row rescaling $L|_{D_{f,*}}=(-1)^fD$ therefore obeys the left
Leibniz rule for $\star$, and $L^2=-D^2=0$.  Row rescaling preserves
every kernel, image, and line of possible cycles.  A further common
unit puts $L(c')=ab^3$ as in Section~\ref{sec:finite}.
The permitted unit in the $d_9$ formula absorbs the analogous convention
change there.
\end{proof}

\subsection{Elimination and the literal affine systems}

We describe the elimination convention so the recipe specifies finite
scalar operations throughout.  Choose any fixed total order of the
monomials and generators extending their weights; for example order
generators by $(w(g),C(g),(e_\nu))$ and then coefficient multi-indices
lexicographically.  Order slots by $(d,f)$ and matrix entries by row
then column.  Given a sequence of row vectors, maintain a dictionary
of rows with distinct first nonzero positions.  For a row $z$, let $i$
be its first nonzero position and $c=z_i$.  If a stored row $p_i$ has
pivot one at $i$, replace $z$ by $z-cp_i$ and repeat at the new first
nonzero position until the row is zero or has an unused pivot.
Discard a zero row.  At an unused pivot, store $c z$ as
$p_i$ and continue with the next input row.  The normalization uses
$c^{-1}=c$ for $c=1,2$ in $\F_3$.  The number of stored rows is the
rank.  Recording the same operations on an identity matrix gives
witnesses for dependencies or for solutions.

To solve a homogeneous system, assign independent parameters to the
nonpivot variables and substitute the pivot equations in reverse
order.  For an affine system append a last constant column and put its
variable equal to one.  A pivot in that last position is an
inconsistency certificate; otherwise the same substitution gives a
constant solution plus a basis of homogeneous solutions.  Kernel,
image, intersection, and quotient computations below use only this
procedure.  The ranks and intrinsic subspaces are independent of the
chosen orders and Ext representatives.

For $n=60$ or $92$, let
\begin{equation}\label{alg:finite-grid}
 \mathcal G_n=\{(f,d):0\leq d\leq n,\quad 0\leq f\leq32,
                                      \quad2d-3f\geq0\}.
\end{equation}
For every pair $T,T+J$ in this grid, $J=(5,1)$, introduce all scalar
entries of a matrix $L_T:D_T\to D_{T+J}$ as independent unknowns.
Slots of dimension zero contribute no entries.  Put $L_T=0$ at
filtration zero.  For each $(k,e)\in\F_3^2$ use the factors
\[
 x=a,v,h,b,c',r,s',y+e t_7'.
\]
For each complete square $T,T+J,T+x,T+x+J$ in the grid, set
\begin{equation}\label{alg:leibniz-equations}
 L_{T+x}\operatorname{mult}(x,T)
 -(-1)^{f_x}\operatorname{mult}(x,T+J)L_T=Z_{x,T},
\end{equation}
where
\begin{equation}\label{alg:leibniz-constants}
 Z_{x,T}=
 \begin{cases}
 \operatorname{mult}(a,T+3b)\operatorname{mult}(b^3,T),&x=c',\\
 k\operatorname{mult}(v,T+3b)\operatorname{mult}(b^3,T),&x=r,\\
 0,&\text{otherwise}.
 \end{cases}
\end{equation}
For $s'$ and $y+e t_7'$ omit the entire equation if either required
factor multiplication has output weight greater than $91$.
Include an equation only when every required matrix is available,
including intermediate matrices in the displayed successive products.
In particular a target lying outside the grid is not silently replaced
by a known zero space.  This is the literal finite system used in
Table~\ref{finite:affine-ranks}.

Every matrix equation here is a specified list of scalar affine
equations.  More explicitly its $(\alpha,\beta)$ entry is
\begin{equation}\label{alg:scalar-leibniz}
 \sum_i (L_{T+x})_{\alpha i}
              \operatorname{mult}(x,T)_{i\beta}
 -(-1)^{f_x}\sum_j
              \operatorname{mult}(x,T+J)_{\alpha j}
              (L_T)_{j\beta}-(Z_{x,T})_{\alpha\beta}=0.
\end{equation}
The named products and coefficients in this formula were fully
specified above; in particular the unknown differential is never used
to supply an Ext product.

For $n=60$, solve these systems for all nine $(k,e)$.  For $n=92$, use
$(k,e)=(1,0)$ and refine its affine solution space as follows.  Express
all entries as affine linear polynomials in its free parameters
$u_1,\ldots,u_a$.  For every complete triple $T,T+J,T+2J$ expand every
entry of
\begin{equation}\label{alg:square-polynomials}
 L_{T+J}(u)L_T(u)
\end{equation}
as a polynomial of degree at most two.  Collect the commutative
monomials $u_i u_j$ with $i\leq j$, then $u_i$, then the constant
monomial.  Eliminate with all quadratic columns placed first.  The
rows remaining with no quadratic term are affine linear equations;
adjoin them to the system, reparameterize, and repeat until no new
restriction is obtained.

\begin{lemma}\label{alg:square-soundness}
Every actual square-zero differential satisfying the imposed Leibniz
equations remains in the affine space at every refinement step.
The refinement terminates.  It makes no assertion that every remaining
point is square-zero.
\end{lemma}

\begin{proof}
Each row produced by elimination is an $\F_3$-linear combination of
the entries of~\eqref{alg:square-polynomials}.  It therefore vanishes
at each actual square-zero differential.  Rows with no quadratic term
are valid affine restrictions on the present parameter space.
Substitution of an affine parametrization preserves all such points.
Every nontrivial refinement decreases the finite affine dimension,
which proves termination.  Retaining only the linear consequences
uses no claim about all polynomial consequences or about the full set
of $\F_3$-points of the quadratic system.
\end{proof}

\subsection{Universal subspaces and the final rank tests}

For an outgoing or incoming matrix obtained by projection from the final
affine space, write its parametrization as
$A(u)=A_0+\sum_{i=1}^a u_iA_i$.  Zero or redundant parameter matrices
may be discarded.  For a column matrix $X$, let $\langle X\rangle$
denote its column span and let $X^\perp$ denote a row matrix whose
kernel is $\langle X\rangle$.  Such a matrix is obtained by computing
$\ker X^{\mathsf t}$ and transposing a basis.  The following formulas
are explicit matrix versions of the enclosures in
Section~\ref{sec:finite}:
\begin{align}
 K^-&=\ker\begin{bmatrix}A_0\\A_1\\\vdots\\A_a\end{bmatrix},
 &K^+&=\ker\bigl([A_1|\cdots|A_a]^\perp A_0\bigr),
 \label{alg:kernel-enclosures}\\
 I^-&=\left\langle A_0\ker\begin{bmatrix}A_1\\\vdots\\A_a\end{bmatrix}
                              \right\rangle,
 &I^+&=\langle[A_0|A_1|\cdots|A_a]\rangle.
 \label{alg:image-enclosures}
\end{align}
When a space appears as a matrix operand, take any column basis of it.
An empty stack of constraints has the full space as its kernel.
Negative-filtration groups are zero.  At other unavailable incoming
or outgoing pairs use the zero space for the lower bound and the
entire relevant space for the upper bound.

\begin{lemma}\label{alg:enclosures-valid}
For every actual matrix $A(u)$ in the affine family,
\[
 K^-\subseteq\ker A(u)\subseteq K^+,
 \qquad I^-\subseteq\im A(u)\subseteq I^+.
\]
These inclusions remain valid when parameters shared with other
matrices are correlated or satisfy further equations.
\end{lemma}

\begin{proof}
Every $A_i$ kills $K^-$.  If $A(u)x=0$, then
$A_0x=-\sum_{i>0}u_iA_ix$ lies in $\sum_{i>0}\im A_i$, which is
the condition defining $K^+$.  On the common kernel of the $A_i$ with
$i>0$, the map $A(u)$ equals $A_0$, proving the lower image inclusion.
The upper image inclusion follows by summing the images of its
summands.  Each argument applies to every single parameter value, so
restricting the allowed parameter values cannot invalidate it.
\end{proof}

Use outgoing matrices at a slot $q$ for $K_q^\pm$ and incoming matrices
for $I_q^\pm$.  Here are matrix recipes for every further test; they
also fix precisely which preimages and intersections are meant.
For the two sum dimensions in Section~\ref{sec:finite}, compute the
ranks of
\[
 [\operatorname{mult}(v,(10,82))K^+_{10,82}\mid I^-_{10,83}],
 \qquad[K^+_{22,86}\mid I^-_{22,86}].
\]
For the cancellation test with $m,h_0,\ell,u,p,z$ as specified there, form
\begin{align}
 W_p&=[I^-_p\mid\operatorname{mult}(b^\ell,z)K^-_z],\label{alg:allowed-space}\\
 A&=K^-_u\ker\bigl(W_p^\perp
                         \operatorname{mult}(P,u)K^-_u\bigr),
                         \label{alg:cancellation-preimage}\\
 F&=[I^-_m\mid\operatorname{mult}(b^{h_0},u)A].
                         \label{alg:cancellation-span}
\end{align}
Take column spans in the second and third lines when computing
dimensions.  The tested triple is
$\rank A,\rank F,\rank[F\mid K^+_m]$.
For an injection test at $q=(f,d)$ with exponent $n$, compute
\begin{equation}\label{alg:injection-test}
 \rank\left[I^-_q\ \middle|\quad
 K^+_q\ker\left((I^+_{q+nb})^\perp
                    \operatorname{mult}(b^n,q)K^+_q\right)\right].
\end{equation}
Together with $\dim I^-_q$, this is exactly the pair recorded in
Table~\ref{finite:injection-data}.  One can avoid annihilator matrices
altogether: the preimage of $\langle H\rangle$ under a matrix $M$ is
the projection to the source of $\ker[M\mid-H]$.  This gives an
equivalent elimination recipe for all the preimages above.

There is a shorter calculation for the second cancellation triple.
Here
\[
 u=(12,76),\quad p=(15,92),\quad z=(9,83),\quad
 W_p=I^-_{15,92}+b^3K^-_{9,83}.
\]
Compute the image of $a:D_{14,91}\to D_{15,92}$.  The ranks of
$W_p$, that image, and their sum are all three.  Hence
\begin{equation}\label{alg:second-factorization}
 \im(a:D_{14,91}\to D_{15,92})\subseteq W_p.
\end{equation}
Since $P=aB(t_4')$, every $Px$ in this slot lies in this image.
Equation~\eqref{alg:cancellation-preimage} therefore has
$A=K^-_{12,76}$, of dimension six.  Only the $a$ matrix and the
$b^3$ matrix are required for this containment calculation.  The
first cancellation triple uses the actual $P$ matrix from $(4,69)$
to $(7,85)$.

\subsection{Arithmetic checkpoints and reproducibility}

\begin{proposition}\label{alg:table-verification}
The recurrences and elimination instructions in this appendix give
the low-cohomology ranks, the affine ranks, the sixteen enclosure rows,
the sum dimensions, the two cancellation triples, the three injection
tests, and the five ordinary product ranks displayed in
Section~\ref{sec:finite}.  In particular the weight-$60$ system is
consistent exactly for $(k,e)=(1,0)$, the stable weight-$92$ affine
space has dimension $226$, and $B(y)^3=0$ in $D_{6,81}$.
\end{proposition}

\begin{proof}[Finite calculation]
The following intermediate outputs specify checks on the elimination
at each stage, in addition to the final tables.  They can be
reconstructed without any external list of Ext groups.

First enumerate~\eqref{alg:generators} and~\eqref{alg:cochain-basis},
construct~\eqref{alg:deformation}, and form~\eqref{alg:cochain-matrix}.
For $(f,d)=(1,27),(6,10),(6,81)$ the dimensions of the three
successive cochain terms are respectively
\[
 (12,49,94),\qquad(3,2,1),\qquad(9301,11419,12864).
\]
Reducing the incoming and outgoing columns gives ranks
\[
 (11,35),\qquad(1,0),\qquad(5256,6156).
\]
For the first slot append, in order, the cochains
$v^{26}a,y,t_7'$; the rank becomes $14$.  For the second append
$vb^3$; the rank becomes $2$.  For the third construct $B(y)$ by
\eqref{alg:bockstein}, apply~\eqref{alg:get-two}, and form its cube
using~\eqref{alg:lift-recursion}--\eqref{alg:product}.  Appending
this cubic leaves the incoming rank $5256$.  Keeping witnesses in
column elimination expresses the cubic as a boundary.  In these
coordinates $B(y)$ has eight coefficient terms on six generators,
the cubic has $76$ terms, and the resulting preimage can be taken to
have $2642$ terms.  Substitution in~\eqref{alg:cochain-matrix} verifies
its boundary directly.  This proves the algebraic zero in the statement.

For weight $60$, the grid contains $1213$ slots and its complete pairs
have $2779$ scalar unknowns.  Expanding the zero and Leibniz equations
gives $15852$ scalar equations.  For each of the nine $(k,e)$ choices,
the coefficient rank is $2667$.  The augmented rank is $2667$ for
$(1,0)$ and $2668$ for each of the other eight choices.  These values
are obtained from~\eqref{alg:scalar-leibniz}, including its constant
column; the pivot at the constant position is an explicit linear
combination giving $1=0$ in each inconsistent case.

At weight $92$ there are $2269$ slots, of total Ext dimension $6341$,
and $1820$ differential pairs with $23384$ entries.  The Leibniz
assembly uses $15453$ factor multiplication matrices, including
successive factors in~\eqref{alg:leibniz-constants}.  There are
$137471$ scalar zero and Leibniz equations, $110060$ of them having
nonzero augmented row.  Coefficient and augmented ranks are both
$23119$.  Thus the initial affine dimension is $265$.  The successive
square eliminations have the following checkpoints:
\[
\begin{array}{r|r|r|r}
 \text{parameters}&\text{scalar square equations}&
       \text{polynomial row rank}&\text{new affine restrictions}\\\hline
 265&16412&72&38\\
 227&16412&20&1\\
 226&16412&19&0
\end{array}
\]
In each row quadratic monomials precede linear ones, as specified
after~\eqref{alg:square-polynomials}; after the final refinement the affine
rank is $23158$.  Apply~\eqref{alg:kernel-enclosures} and
\eqref{alg:image-enclosures} to its projected matrices.  They give
the sixteen rows of Table~\ref{finite:enclosures}.  The two additional
stacked ranks are $1$ and $2$.  Equations
\eqref{alg:allowed-space}--\eqref{alg:cancellation-span}, using
\eqref{alg:second-factorization} in the second case, give
$(4,5,5)$ and $(6,5,5)$.  Equation~\eqref{alg:injection-test} gives
the pairs $(3,3),(2,2),(4,4)$.

Finally compute the five iterated $b$ maps of
Table~\ref{finite:product-ranks} with cutoff $93$.  As a check
independent of quotient-coordinate extraction, append the resulting
product cochains to the incoming coboundary columns in their target
slots.  In the order of that table, incoming ranks and ranks after
appending are
\[
 (4491,4499),\quad(7540,7546),\quad(7589,7595),\quad
 (1868,1872),\quad(1000,1004).
\]
Their differences are the product ranks $8,6,6,4,4$.
All these operations have been reduced to the polynomial recursions,
the terminating lifting recursion, and the explicitly stated
elimination rules.  The equalities used to infer the spectral-sequence
conclusions are therefore finite algebra calculations; their use is
proved in Section~\ref{sec:finite}.
\end{proof}

\begin{remark}[Distinct forms of computational checking]
Reconstructing the larger tables by another implementation of this
same recurrence checks its arithmetic, while sharing the mathematical
resolution construction.  There are also ordinary normalized-cobar
calculations which do not use the smaller resolution.  For example
in weight $27$ their dimensions in filtrations zero and one are $12$
and $470$, and their incoming and outgoing ranks at filtration one
are $11$ and $456$.  They give the same three-dimensional group with
the representatives $v^{26}a,y,t_7'$.  At $(6,10)$ their middle
dimension is $253$, the incoming and outgoing ranks are $162$ and
$90$, and appending $vb^3$ increases the incoming rank by one.
These different complex sizes distinguish the ordinary-cobar checks
from a replay of the smaller-resolution calculation.  Verification of
the cubic boundary, substitution into the affine equations, and dense
elimination for the final subspaces provide further arithmetic
checks.  None replaces the resolution and comparison proof above or
the topological arguments of Section~\ref{sec:finite}.
\end{remark}
